%%%%%%%%%%%%%%%%%%%%%%%%%%%%%%%%%%%%%%%%%%%%%%%%%%%%%%%%%%%%%%
%% Web GIS Laboratory 4: Getting started with GeoDjango
%% Institute of Geographical Information Systems, NUST
%%
%% Layout follows "A Psychology Lab Manual" by John Paul Minda
%% (Overleaf, CC BY 4.0): article class, fullpage margins,
%% Helvetica throughout, parskip, blue links, no colour boxes.
%%
%% COMPILER: pdfLaTeX
%%%%%%%%%%%%%%%%%%%%%%%%%%%%%%%%%%%%%%%%%%%%%%%%%%%%%%%%%%%%%%

\documentclass{article}
\usepackage{fullpage}

\renewcommand{\familydefault}{\sfdefault}
\usepackage[scaled=1]{helvet}
\usepackage[helvet]{sfmath}
\everymath={\sf}

\usepackage{parskip}
\usepackage[colorlinks=true, allcolors=blue]{hyperref}

%% Required by this document but not by the original template.
%% Encoding first: without T1, < and > inside code print as
%% inverted punctuation and straight quotes turn curly.
\usepackage{iftex}
\ifPDFTeX
  \usepackage[T1]{fontenc}
  \usepackage[utf8]{inputenc}
\fi
\usepackage{textcomp}
\usepackage{upquote}
\usepackage{xcolor}
\usepackage{listings}
\usepackage{booktabs}
\usepackage{tabularx}
\usepackage{array}
\usepackage{enumitem}

\setlist[enumerate]{itemsep=.25em,topsep=.5em}
\setlist[itemize]{itemsep=.25em,topsep=.5em}
\renewcommand{\arraystretch}{1.25}
\newcolumntype{L}[1]{>{\raggedright\arraybackslash}p{#1}}

%% Inline code. \cd for short tokens, \cdb for long paths and URLs,
%% which must be allowed to break across lines.
\urlstyle{tt}
\newcommand{\cd}[1]{\texttt{#1}}
\newcommand{\cdb}[1]{\nolinkurl{#1}}

%% Code blocks: a plain rule, no shading, in keeping with the template
\lstdefinestyle{plain}{
  basicstyle=\ttfamily\small,
  frame=single,
  rulecolor=\color{black!35},
  framesep=7pt,
  xleftmargin=7pt,
  xrightmargin=7pt,
  breaklines=true,
  columns=fullflexible,
  keepspaces=true,
  showstringspaces=false,
  upquote=true,
  aboveskip=1em,
  belowskip=1em,
}
\lstset{style=plain}
\lstnewenvironment{code}{\lstset{style=plain}}{}

%% Caption for a code block. Deliberately not \paragraph, which is a
%% run-in heading and would be overlapped by the listing frame.
\newcommand{\codecap}[1]{%
  \par\vspace{.9em}\noindent\textbf{#1}\par\vspace{-.6em}}

%% A block that must not be skimmed past: ruled above and below,
%% with a bold label. Used for warnings and checkpoints.
\newenvironment{keypoint}[1]
  {\par\vspace{.7em}\noindent\rule{\linewidth}{.6pt}\par\vspace{.15em}
   \noindent\textbf{#1}\par\nobreak\vspace{-.3em}}
  {\par\vspace{.15em}\noindent\rule{\linewidth}{.6pt}\par\vspace{.7em}}

\title{Web GIS Laboratory 4\\[.3em]\large Getting started with GeoDjango:
       Django, PostGIS, the admin panel and Leaflet}
\author{Dr.\ Shahid Nawaz Khan\\
        Institute of Geographical Information Systems\\
        School of Civil and Environmental Engineering\\
        National University of Sciences and Technology, Islamabad}
\date{MS Remote Sensing and GIS}

\setcounter{tocdepth}{1}

\begin{document}
\maketitle

\begin{tabularx}{\textwidth}{@{}L{3.4cm} X@{}}
\toprule
Duration    & 3 hours \\
Builds on   & Labs 2 and 3: the same districts, now in a database \\
Software    & Python, Django, PostgreSQL with PostGIS, OSGeo4W, VS Code \\
Deliverable & A working GeoDjango project in your repository, and a report \\
\bottomrule
\end{tabularx}

\tableofcontents


%==============================================================
\section*{Before you begin}
\addcontentsline{toc}{section}{Before you begin}

\subsection*{What changes in this lab}

Until now the browser did everything. It fetched a GeoJSON file that sat
next to the page and drew it. That works for a fixed dataset, but nothing
can be edited, nothing is checked, and every user downloads everything.

From this lab on, the data lives in a \textbf{spatial database}, and a
\textbf{web framework} stands between the database and the browser:

\vspace{.4em}
\begin{center}
\begin{tabular}{c c c c c}
\textbf{PostGIS} & $\longleftrightarrow$ & \textbf{Django} &
$\longrightarrow$ & \textbf{Leaflet in the browser} \\[.2em]
\small stores and queries & & \small decides what to send & &
\small draws what it receives \\
\end{tabular}
\end{center}
\vspace{.4em}

GeoDjango is the part of Django that understands geometry. By the end
of the lab the 160 districts from Labs 2 and 3 will be stored in PostGIS,
editable through Django's admin panel, and drawn by Leaflet from data
Django serves.

\subsection*{Learning outcomes}

\begin{enumerate}
  \item Create a Python virtual environment and install Django into it.
  \item Install PostgreSQL with the PostGIS extension, and create a
        spatial database.
  \item Install the GDAL, GEOS and PROJ libraries GeoDjango depends on,
        and connect Django to PostGIS.
  \item Define a model with a geometry field, and read the SQL Django
        generates for it.
  \item Load a GeoJSON file into the database.
  \item Use the Django admin panel to browse, filter, search and edit
        spatial data.
  \item Serve data from the database as GeoJSON and draw it with Leaflet.
\end{enumerate}

\subsection*{Structure and timing}

This lab is mostly installation. Installation fails in different ways on
different machines, so the times are a guide. Parts 1 to 7 are the core;
Part 8 is an extension.

\vspace{.6em}
\begin{tabularx}{\textwidth}{@{}c X r@{}}
\toprule
\textbf{Part} & \textbf{What you do} & \textbf{Time} \\
\midrule
1 & Python and a first Django project            & 20 min \\
2 & Install PostgreSQL and PostGIS               & 25 min \\
3 & Connect Django to PostGIS                    & 25 min \\
4 & A model with a geometry                      & 15 min \\
5 & Load the districts                           & 15 min \\
6 & The admin panel                              & 20 min \\
7 & Leaflet, fed by Django                       & 25 min \\
\midrule
8 & Ask PostGIS a spatial question \textit{(extension)} & 15 min \\
\midrule
9 & Commit and submit                            & 10 min \\
\bottomrule
\end{tabularx}

\subsection*{What you need}

\begin{tabularx}{\textwidth}{@{}L{4.2cm} X@{}}
\toprule
\textbf{Item} & \textbf{Requirement} \\
\midrule
Operating system & Windows 10 or 11 \\
Disk space       & About 2 GB for PostgreSQL, PostGIS and OSGeo4W \\
Rights           & Administrator rights, for the two installers \\
Data             & \cd{districts.geojson} from Lab 3, with \cd{pop\_2023} \\
Internet         & The installers download several hundred megabytes \\
\bottomrule
\end{tabularx}

\begin{keypoint}{Keep a record as you go}
Write down every version number you see, and copy any error message in
full before you try to fix it. Part of your report is a record of what
went wrong on your machine and how you fixed it.
\end{keypoint}


%==============================================================
\section{Python and a first Django project}

\subsection{Install Python}

Django 6.1 needs Python 3.12 or later. Check whether you already have it:

\begin{code}
python --version
\end{code}

If the command is not found, or prints a version below 3.12, install
Python 3.13 from \cd{python.org}. On the first screen of the installer,
tick \textbf{Add python.exe to PATH} before clicking Install. Close
PowerShell, open a new one, and check the version again.

\subsection{Create the project folder and a virtual environment}

Work inside your course repository, in a new folder:

\begin{code}
cd C:\webgis-lab
mkdir lab4
cd lab4
python -m venv venv
venv\Scripts\Activate.ps1
\end{code}

The prompt now starts with \cd{(venv)}. Everything you install with
\cd{pip} from here goes into the \cd{venv} folder, not into Python
itself.

\begin{keypoint}{If PowerShell refuses to run Activate.ps1}
A message saying that running scripts is disabled on this system means
PowerShell's execution policy blocks it. Allow scripts for your own
account, once, then activate again:

\begin{code}
Set-ExecutionPolicy -Scope CurrentUser RemoteSigned
\end{code}
\end{keypoint}

\paragraph{Why a virtual environment}
Each project gets its own copy of its libraries, at its own versions. A
project from another course that needs a different Django cannot break
this one, and \cd{requirements.txt} records exactly what this project
needs.

\subsection{Install Django and the PostgreSQL driver}

\begin{code}
pip install django "psycopg[binary]"
pip freeze > requirements.txt
python -m django --version
\end{code}

\cd{psycopg} is how Python talks to PostgreSQL. The \cd{[binary]} part
installs a ready-built version, so no compiler is needed.

\subsection{Start a project and an app}

\begin{code}
django-admin startproject webgis .
python manage.py startapp districts
\end{code}

Note the dot at the end of the first line: it puts the project in the
current folder instead of a new one.

\begin{tabularx}{\textwidth}{@{}L{4.4cm} X@{}}
\toprule
\textbf{File or folder} & \textbf{Holds} \\
\midrule
\cd{manage.py}            & The command you run everything through \\
\cd{webgis/settings.py}   & The project's configuration \\
\cd{webgis/urls.py}       & Which URL goes to which code \\
\cd{districts/}           & An \textbf{app}: one part of the project, here
                            everything about districts \\
\bottomrule
\end{tabularx}

\subsection{Run it}

\begin{code}
python manage.py migrate
python manage.py runserver
\end{code}

Open \cdb{http://127.0.0.1:8000/}. You should see a page saying the
install worked. Then open \cdb{http://127.0.0.1:8000/admin/}: a login page.
There is no user yet; that comes in Part 6.

\cd{migrate} created a database. By default it is SQLite, a single file
called \cd{db.sqlite3}. SQLite cannot store geometry properly, so Part 3
replaces it. Stop the server with \cd{Ctrl+C} and delete
\cd{db.sqlite3}.

\begin{keypoint}{Checkpoint 1}
A screenshot of the ``install worked'' page with the address bar
visible, and the output of \cd{python -m django --version}.
\end{keypoint}


%==============================================================
\section{Install PostgreSQL and PostGIS}

\subsection{PostgreSQL}

From \cd{postgresql.org}, under Download, Windows, take the installer
for the latest version (17 or later). Run it, and accept the defaults
except for these:

\begin{tabularx}{\textwidth}{@{}L{4.4cm} X@{}}
\toprule
\textbf{Screen} & \textbf{What to do} \\
\midrule
Select Components & Keep all four: the server, pgAdmin 4, Stack Builder
                   and the command line tools \\
Password          & The password for the \cd{postgres} superuser. Choose
                    one and \textbf{write it down}; it cannot be
                    recovered \\
Port              & Leave it at \cd{5432} \\
Last screen       & Leave \textbf{Launch Stack Builder} ticked, and
                    click Finish \\
\bottomrule
\end{tabularx}

\subsection{PostGIS, through Stack Builder}

PostgreSQL on its own knows nothing about geometry. PostGIS is an
extension that adds geometry types, spatial indexes and several hundred
spatial functions.

\begin{enumerate}
  \item In Stack Builder, choose your PostgreSQL installation from the
        list, then Next.
  \item Open \textbf{Spatial Extensions} and tick the \textbf{PostGIS}
        bundle for your PostgreSQL version.
  \item Download, install, and accept the defaults throughout, including
        any questions about environment variables.
\end{enumerate}

\subsection{Create a spatial database}

Open \textbf{pgAdmin 4}, and connect to the server with your password.

\begin{enumerate}
  \item Right-click \textbf{Databases}, \textbf{Create},
        \textbf{Database}. Name it \cd{webgis}, and Save.
  \item Select \cd{webgis}, then \textbf{Tools}, \textbf{Query Tool}.
  \item Run these two statements:
\end{enumerate}

\begin{code}
CREATE EXTENSION postgis;
SELECT postgis_full_version();
\end{code}

The second returns one long line listing PostGIS, GEOS and PROJ with their
versions. That line is your proof that the database can hold geometry.

\begin{keypoint}{Checkpoint 2}
A screenshot of pgAdmin with the \cd{webgis} database selected and the
result of \cd{postgis\_full\_version()} visible.
\end{keypoint}


%==============================================================
\section{Connect Django to PostGIS}

Three things are needed: the libraries GeoDjango uses to handle geometry
in Python, the password, and the settings.

\subsection{GDAL, GEOS and PROJ, through OSGeo4W}

GeoDjango does not do the geometry work itself. It calls three C
libraries: \textbf{GDAL} to read and write formats, \textbf{GEOS} for
geometry operations, \textbf{PROJ} for coordinate systems. QGIS uses the
same three. On Windows, the usual way to get them is OSGeo4W.

\begin{enumerate}
  \item Download \cd{osgeo4w-setup.exe} from \cd{trac.osgeo.org/osgeo4w}.
  \item Run it and choose \textbf{Express Install}, then the default
        download site.
  \item From the list of packages, tick \textbf{GDAL} only, and continue.
        It installs to \cd{C:\textbackslash OSGeo4W}.
\end{enumerate}

\begin{keypoint}{If the OSGeo4W installer will not run}
QGIS for Windows is built on OSGeo4W and carries the same libraries. If
you installed QGIS for Lab 2, you can point Django at its folder instead,
for example \cd{C:\textbackslash Program Files\textbackslash QGIS 3.40.10}.
Use the exact folder name on your machine in the settings below. Use this
only if OSGeo4W fails.
\end{keypoint}

\subsection{Store the password outside the code}

Your settings file will be committed to GitHub, so the database password
must not be written in it. Store it in an environment variable instead:

\begin{code}
setx WEBGIS_DB_PASSWORD "the-password-you-chose"
\end{code}

\cd{setx} saves the variable for every \textbf{new} window. Close VS Code
completely and open it again, then activate the environment again with
\cd{venv\textbackslash Scripts\textbackslash Activate.ps1}.

\subsection{Edit the settings}

Open \cd{webgis/settings.py} and make three changes.

\paragraph{One: tell Python where OSGeo4W is}
Replace the line \cd{from pathlib import Path} near the top, and add the
block below \cd{BASE\_DIR}:

\codecap{webgis/settings.py, top of the file}
\begin{code}
import os
from pathlib import Path

# Build paths inside the project like this: BASE_DIR / 'subdir'.
BASE_DIR = Path(__file__).resolve().parent.parent

# GeoDjango needs the GDAL, GEOS and PROJ libraries. On Windows they come
# from OSGeo4W; putting its bin folder on PATH lets Django find them.
if os.name == "nt":
    OSGEO4W = r"C:\OSGeo4W"
    os.environ["OSGEO4W_ROOT"] = OSGEO4W
    os.environ["GDAL_DATA"] = OSGEO4W + r"\share\gdal"
    os.environ["PROJ_DATA"] = OSGEO4W + r"\share\proj"
    os.environ["PROJ_LIB"] = OSGEO4W + r"\share\proj"
    os.environ["PATH"] = OSGEO4W + r"\bin;" + os.environ["PATH"]
\end{code}

\cd{os.name == "nt"} is true only on Windows, so the same settings file
still works on a Linux server later.

\paragraph{Two: add GeoDjango and your app}

\codecap{webgis/settings.py, INSTALLED\_APPS}
\begin{code}
INSTALLED_APPS = [
    "django.contrib.admin",
    "django.contrib.auth",
    "django.contrib.contenttypes",
    "django.contrib.sessions",
    "django.contrib.messages",
    "django.contrib.staticfiles",
    "django.contrib.gis",
    "districts",
]
\end{code}

\paragraph{Three: replace SQLite with PostGIS}

\codecap{webgis/settings.py, DATABASES}
\begin{code}
# The password is read from an environment variable so that it is never
# written into a file that gets committed to GitHub.
DATABASES = {
    "default": {
        "ENGINE": "django.contrib.gis.db.backends.postgis",
        "NAME": "webgis",
        "USER": "postgres",
        "PASSWORD": os.environ.get("WEBGIS_DB_PASSWORD", ""),
        "HOST": "localhost",
        "PORT": "5432",
    }
}
\end{code}

\subsection{Check each piece separately}

Run these one at a time. Each tests one link in the chain, so the first
one to fail tells you where the problem is.

\begin{code}
python manage.py shell -c "from django.contrib.gis import gdal, geos; print(gdal.gdal_version(), geos.geos_version())"
python manage.py check
python manage.py migrate
\end{code}

\begin{tabularx}{\textwidth}{@{}L{3.2cm} X@{}}
\toprule
\textbf{Command} & \textbf{Proves} \\
\midrule
The \cd{shell} line & Python can load GDAL and GEOS. It prints their
                      version numbers \\
\cd{check}          & The settings are valid \\
\cd{migrate}        & Django can log in to PostGIS and create its tables
                      there \\
\bottomrule
\end{tabularx}

\vspace{.8em}
In pgAdmin, refresh \cd{webgis}, then open \textbf{Schemas},
\textbf{public}, \textbf{Tables}. Django's own tables are now there,
next to \cd{spatial\_ref\_sys}, which belongs to PostGIS.

\begin{keypoint}{Checkpoint 3}
The output of the three commands, and a screenshot of the tables in
pgAdmin.
\end{keypoint}


%==============================================================
\section{A model with a geometry}

\subsection{The model}

A \textbf{model} is a Python class that Django turns into a database
table. Each attribute becomes a column.

\codecap{districts/models.py, complete}
\begin{code}
from django.contrib.gis.db import models


class District(models.Model):
    name = models.CharField(max_length=100)
    province = models.CharField(max_length=50)
    pcode = models.CharField(max_length=10, unique=True)
    pop_2023 = models.IntegerField(null=True, blank=True)
    area_km2 = models.FloatField()
    geom = models.MultiPolygonField(srid=4326)

    class Meta:
        ordering = ["province", "name"]

    def __str__(self):
        return f"{self.name} ({self.province})"
\end{code}

\begin{tabularx}{\textwidth}{@{}L{4.4cm} X@{}}
\toprule
\textbf{Line} & \textbf{Why} \\
\midrule
\cd{from django.contrib.gis.db} & The GeoDjango version of \cd{models},
                                  which adds the geometry fields \\
\cd{null=True, blank=True}      & 25 districts have no census figure.
                                  Empty must be allowed \\
\cd{MultiPolygonField}          & Some districts are in more than one
                                  piece \\
\cd{srid=4326}                  & WGS 84 longitude and latitude, the same
                                  as the GeoJSON \\
\cd{\_\_str\_\_}                & How a district is named in the admin and
                                  in the shell \\
\bottomrule
\end{tabularx}

\subsection{From model to table}

\begin{code}
python manage.py makemigrations districts
python manage.py sqlmigrate districts 0001
python manage.py migrate
\end{code}

\cd{makemigrations} writes a migration file describing the change.
\cd{sqlmigrate} shows the SQL that file will run, without running it.
Read it. Two lines matter:

\begin{code}
"geom" geometry(MULTIPOLYGON,4326) NOT NULL
CREATE INDEX "districts_district_geom_..." ON "districts_district" USING GIST ("geom");
\end{code}

The first is a real PostGIS geometry column, constrained to
multipolygons in EPSG:4326. The second is a \textbf{spatial index}. You
did not ask for it; GeoDjango adds one to every geometry field, because
without it every spatial query would have to examine every row.

\begin{keypoint}{Checkpoint 4}
The output of \cd{sqlmigrate}, and a screenshot of the
\cd{districts\_district} table's columns in pgAdmin.
\end{keypoint}


%==============================================================
\section{Load the districts}

\subsection{The data file}

\begin{code}
mkdir data
copy ..\lab3\data\districts.geojson data\districts.geojson
\end{code}

\subsection{A loading script}

\codecap{districts/load.py, complete}
\begin{code}
from pathlib import Path

from django.contrib.gis.gdal import DataSource
from django.contrib.gis.geos import MultiPolygon

from .models import District

data_file = Path(__file__).resolve().parent.parent / "data" / "districts.geojson"


def run():
    layer = DataSource(data_file)[0]
    print(layer.name, ":", len(layer), "features,", layer.geom_type, "geometry")

    # Start from an empty table, so the script can be run more than once
    District.objects.all().delete()

    for feature in layer:
        geom = feature.geom.geos
        geom.srid = 4326

        # The file mixes Polygon and MultiPolygon. The model stores only
        # MultiPolygon, so wrap each single polygon in a multipolygon.
        if geom.geom_type == "Polygon":
            geom = MultiPolygon(geom, srid=4326)

        District.objects.create(
            name=feature.get("district"),
            province=feature.get("province"),
            pcode=feature.get("pcode"),
            pop_2023=feature.get("pop_2023"),
            area_km2=feature.get("area_km2"),
            geom=geom,
        )

    print(District.objects.count(), "districts saved")
\end{code}

Run it:

\begin{code}
python manage.py shell -c "from districts import load; load.run()"
\end{code}

It should report 160 features and 160 districts saved. Note the geometry
type it prints for the layer.

\paragraph{Why the Polygon check}
The file holds 152 \cd{Polygon} features and 8 \cd{MultiPolygon}
features, so GDAL reports the layer's geometry type as \cd{Unknown}. The
column accepts only multipolygons, and PostGIS would reject a plain
polygon. Wrapping it in a \cd{MultiPolygon} of one part changes nothing
on the map and satisfies the column. GeoDjango has a shortcut for loading
files, \cd{LayerMapping}, but it refuses a layer of mixed type outright,
which is why this script reads the features one by one.

\subsection{Talk to the data}

Open the Django shell:

\begin{code}
python manage.py shell
\end{code}

and type these one at a time. This is the \textbf{ORM}: Python that
Django turns into SQL.

\begin{code}
District.objects.count()
District.objects.filter(province="Sindh").count()
District.objects.filter(pop_2023__isnull=True).count()
lahore = District.objects.get(name="Lahore")
lahore.pop_2023
lahore.geom.geom_type
lahore.geom.num_coords
exit()
\end{code}

The shell imports your models automatically when it starts, which is why
\cd{District} needs no import.

\begin{keypoint}{Checkpoint 5}
The output of the loading script and of the seven shell lines.
\end{keypoint}


%==============================================================
\section{The admin panel}

Django builds a complete editing interface for your models: login,
lists, search, filters and forms. You only describe what you want to see.

\subsection{Create an administrator}

\begin{code}
python manage.py createsuperuser
\end{code}

Choose a username, an email address and a password. This account is for
the admin panel only; it is not your PostgreSQL password.

\subsection{Register the model}

\codecap{districts/admin.py, complete}
\begin{code}
from django.contrib.gis import admin

from .models import District


@admin.register(District)
class DistrictAdmin(admin.GISModelAdmin):
    list_display = ["name", "province", "pop_2023", "area_km2"]
    list_filter = ["province"]
    search_fields = ["name", "pcode"]
    ordering = ["province", "name"]
\end{code}

\begin{tabularx}{\textwidth}{@{}L{3.6cm} X@{}}
\toprule
\textbf{Option} & \textbf{Gives you} \\
\midrule
\cd{GISModelAdmin}   & A map for the geometry field, instead of a box of
                       text \\
\cd{list\_display}   & The columns of the list page \\
\cd{list\_filter}    & A filter panel on the right, one entry per
                       province \\
\cd{search\_fields}  & A search box that matches name or P-code \\
\bottomrule
\end{tabularx}

\subsection{Explore it}

Run the server and log in at \cdb{http://127.0.0.1:8000/admin/}.

\begin{enumerate}
  \item The home page lists \textbf{Users} and \textbf{Groups}, which
        Django provides, and \textbf{Districts}, which you registered.
  \item Open Districts. Filter by Sindh: the count above the table reads
        29. Search for \cd{karachi}.
  \item Click a column header to sort. Sort by population.
  \item Open Lahore. The geometry appears on an OpenStreetMap basemap,
        drawn by OpenLayers. Zoom in on it.
  \item Change Lahore's population to any other number and Save. In
        pgAdmin, run the query below and confirm that the database has
        changed. Then set the figure back to 13004135.
  \item Click \textbf{History} on Lahore's page to see who changed what.
\end{enumerate}

\begin{code}
SELECT name, pop_2023 FROM districts_district WHERE name = 'Lahore';
\end{code}

\paragraph{What just happened}
The admin did not edit a copy. Every save is an \cd{UPDATE} on the
PostGIS table, and the change history is itself stored in a table that
\cd{migrate} created. Editing the geometry on the map works the same
way: the shape is sent back as geometry and written to the \cd{geom}
column.

\begin{keypoint}{The admin is for people you trust}
Anyone who can log in can change or delete every row. It is a tool for
the people who maintain the data, not a page for the public. In
\textbf{Users} you can create accounts with fewer permissions, for
example one that can change districts but not delete them.
\end{keypoint}

\begin{keypoint}{Checkpoint 6}
Three screenshots: the district list filtered to one province; Lahore's
page with its map; and the pgAdmin query showing your changed
population.
\end{keypoint}


%==============================================================
\section{Leaflet, fed by Django}

In Labs 2 and 3, Leaflet fetched \cd{data/districts.geojson}, a file.
Now it fetches a URL, and Django builds the answer from the database each
time it is asked.

\subsection{Two views}

A \textbf{view} is a Python function that receives a request and returns
a response.

\codecap{districts/views.py}
\begin{code}
from django.core.serializers import serialize
from django.http import HttpResponse
from django.shortcuts import render

from .models import District


def map_page(request):
    return render(request, "districts/map.html")


def districts_geojson(request):
    districts = District.objects.all()

    # /api/districts/?province=Sindh returns one province only
    province = request.GET.get("province")
    if province:
        districts = districts.filter(province=province)

    data = serialize(
        "geojson",
        districts,
        geometry_field="geom",
        fields=["name", "province", "pop_2023", "area_km2"],
    )
    return HttpResponse(data, content_type="application/geo+json")
\end{code}

\cd{serialize("geojson", ...)} turns rows into a GeoJSON
FeatureCollection: the geometry field becomes each feature's
\cd{geometry}, the listed fields become its \cd{properties}.

\subsection{Two URLs}

Create a new file in the app:

\codecap{districts/urls.py, new file}
\begin{code}
from django.urls import path

from . import views

urlpatterns = [
    path("", views.map_page, name="map"),
    path("api/districts/", views.districts_geojson, name="districts-geojson"),
]
\end{code}

and include it from the project:

\codecap{webgis/urls.py}
\begin{code}
from django.contrib import admin
from django.urls import include, path

urlpatterns = [
    path("admin/", admin.site.urls),
    path("", include("districts.urls")),
]
\end{code}

Run the server and open \cdb{http://127.0.0.1:8000/api/districts/}. The
browser shows GeoJSON, built from PostGIS. Then try
\cdb{http://127.0.0.1:8000/api/districts/?province=Sindh} and compare
the number of features.

\subsection{The map page}

Django looks for templates in a \cd{templates} folder inside each app.
Create the folders \cdb{districts/templates/districts/}, and in the
innermost one the file \cd{map.html}.

\codecap{districts/templates/districts/map.html}
\begin{code}
<!DOCTYPE html>
<html lang="en">
<head>
  <meta charset="utf-8">
  <meta name="viewport" content="width=device-width, initial-scale=1">
  <title>Districts from PostGIS</title>
  <link rel="stylesheet"
        href="https://unpkg.com/leaflet@1.9.4/dist/leaflet.css">
  <style>
    body { font-family: system-ui, sans-serif; margin: 0;
           color: #23373b; background: #f4f6f6; }
    header { padding: 16px 24px 12px; }
    h1 { margin: 0; font-size: 22px; }
    .subtitle { margin: 4px 0 0; color: #5b6a6e; }
    #map { height: calc(100vh - 90px); min-height: 400px; }
  </style>
</head>
<body>
  <header>
    <h1>Districts from PostGIS</h1>
    <p class="subtitle">
      <span id="count">loading...</span>, served by Django from the database.
    </p>
  </header>

  <div id="map"></div>

  <script src="https://unpkg.com/leaflet@1.9.4/dist/leaflet.js"></script>
  <script>
    const DISTRICTS_URL = "{% url 'districts-geojson' %}";

    const map = L.map('map', {
      maxBounds: [[21.5, 58.5], [39.0, 80.5]],
      maxBoundsViscosity: 1.0,
      minZoom: 5
    }).setView([30.3753, 69.3451], 5);

    L.tileLayer('https://tile.openstreetmap.org/{z}/{x}/{y}.png', {
      maxZoom: 19,
      attribution: '&copy; OpenStreetMap contributors'
    }).addTo(map);

    const colours = {
      'Punjab': '#1b9e77', 'Sindh': '#d95f02', 'Khyber Pakhtunkhwa': '#7570b3',
      'Balochistan': '#e7298a', 'Islamabad': '#66a61e',
      'Azad Kashmir': '#e6ab02', 'Gilgit Baltistan': '#a6761d'
    };

    function show(value) {
      return value === null || value === undefined
        ? 'not recorded' : value.toLocaleString();
    }

    async function loadDistricts() {
      const response = await fetch(DISTRICTS_URL);
      if (!response.ok) {
        console.error('Request failed with status', response.status);
        return;
      }
      const data = await response.json();

      const layer = L.geoJSON(data, {
        style: function (feature) {
          return {
            color: '#ffffff',
            weight: 1,
            fillColor: colours[feature.properties.province] || '#888888',
            fillOpacity: 0.7
          };
        },
        onEachFeature: function (feature, featureLayer) {
          const p = feature.properties;
          featureLayer.bindTooltip(p.name + ', ' + p.province, { sticky: true });
        }
      }).addTo(map);

      map.fitBounds(layer.getBounds());
      document.getElementById('count').textContent =
        data.features.length + ' districts';
    }

    loadDistricts();
  </script>
</body>
</html>
\end{code}

Open \cdb{http://127.0.0.1:8000/}. Hover over the districts.

\paragraph{\texttt{\{\% url 'districts-geojson' \%\}}}
Django runs the template before the browser sees it. This tag is
replaced with the URL that \cd{urls.py} named \cd{districts-geojson},
so the page reads \cd{const DISTRICTS\_URL = "/api/districts/";}. Press
\cd{Ctrl+U} in the browser to see the page as it arrived. If you ever
move the API to another address, only \cd{urls.py} changes.

\paragraph{Everything you learned still applies}
The JavaScript is the same Leaflet as Labs 2 and 3. Only the source of
the data changed. Change Lahore's province colour in the admin by
editing its province name, reload the map, and see the map follow the
database. Put it back afterwards.

\begin{keypoint}{Checkpoint 7}
A screenshot of the map at \cd{127.0.0.1:8000} with a tooltip showing,
and one of the Network tab with the \cd{api/districts/} request selected
and its size visible.
\end{keypoint}


%==============================================================
\section{Ask PostGIS a spatial question}
\textit{Extension exercise}

So far the database has only stored and returned rows. Its real
strength is answering spatial questions. Here the browser sends a point
and PostGIS says which district contains it.

\subsection{The view}

Add to \cd{districts/views.py}, extending the imports at the top:

\codecap{districts/views.py, added}
\begin{code}
from django.contrib.gis.geos import Point
from django.http import JsonResponse


def district_at(request):
    try:
        lat = float(request.GET["lat"])
        lng = float(request.GET["lng"])
    except (KeyError, ValueError):
        return JsonResponse({"error": "give lat and lng as numbers"}, status=400)

    point = Point(lng, lat, srid=4326)
    district = District.objects.filter(geom__contains=point).first()

    if district is None:
        return JsonResponse({"district": None})
    return JsonResponse({
        "district": district.name,
        "province": district.province,
        "pop_2023": district.pop_2023,
    })
\end{code}

Add its URL to \cd{districts/urls.py}:

\begin{code}
    path("api/where/", views.district_at, name="district-at"),
\end{code}

Test it directly: \cdb{http://127.0.0.1:8000/api/where/?lat=33.6423&lng=72.9906}
should name Islamabad.

\paragraph{\texttt{Point(lng, lat)}}
GEOS, like GeoJSON, puts longitude first. Leaflet, as always, puts
latitude first. The view takes them by name so the order cannot be
confused in the URL.

\paragraph{\texttt{geom\_\_contains}}
This becomes a PostGIS \cd{ST\_Contains} in SQL, and it uses the spatial
index from Part 4 to skip every district whose bounding box does not
include the point.

\subsection{Call it from the map}

In \cd{map.html}, below \cd{const DISTRICTS\_URL}, add
\cd{const WHERE\_URL = "\{\% url 'district-at' \%\}";} and, after
\cd{loadDistricts();}, this handler:

\begin{code}
    map.on('click', async function (e) {
      const url = WHERE_URL + '?lat=' + e.latlng.lat + '&lng=' + e.latlng.lng;
      const answer = await (await fetch(url)).json();
      const text = answer.district
        ? '<strong>' + answer.district + '</strong>, ' + answer.province +
          '<br>Population: ' + show(answer.pop_2023)
        : 'Not inside any district';
      L.popup().setLatLng(e.latlng).setContent(text).openOn(map);
    });
\end{code}

\subsection{The same question in SQL}

In pgAdmin, run the same query by hand, then compare areas:

\begin{code}
SELECT name, province FROM districts_district
WHERE ST_Contains(geom, ST_SetSRID(ST_MakePoint(72.9906, 33.6423), 4326));

SELECT name,
       round(area_km2) AS stored,
       round((ST_Area(geom::geography) / 1e6)::numeric) AS geodesic,
       round((ST_Area(ST_Transform(geom, 3857)) / 1e6)::numeric) AS web_mercator
FROM districts_district
WHERE name IN ('Lahore', 'Gilgit', 'Gwadar');
\end{code}

The second query is Part 9 of Lab 2 done in one statement. Compare the
\cd{web\_mercator} column with the others and relate the ratio to each
district's latitude.

\begin{keypoint}{Checkpoint 8}
A screenshot of a popup produced by clicking the map, and the result of
the area query.
\end{keypoint}


%==============================================================
\section{Commit and submit}

\subsection{What not to commit}

Create \cd{lab4/.gitignore}:

\codecap{lab4/.gitignore}
\begin{code}
venv/
__pycache__/
*.pyc
db.sqlite3
\end{code}

The \cd{venv} folder is hundreds of megabytes and can be rebuilt at any
time from \cd{requirements.txt}. Check with \cd{git status} that it is
not listed before you commit.

\subsection{Commit and push}

\begin{code}
cd C:\webgis-lab
git status
git add -A
git commit -m "Lab 4: GeoDjango with PostGIS, admin and Leaflet"
git push
\end{code}

\begin{keypoint}{This lab has no live URL}
GitHub Pages serves files; it cannot run Python or a database. This
project runs only on your own machine for now. Putting it on a real
server is a later lab. Your submission is the repository and the report.
\end{keypoint}

\subsection{Write the report}

Create \cd{lab4/report.md} containing:

\begin{enumerate}
  \item Your checkpoint screenshots.
  \item The version numbers of Python, Django, PostgreSQL, PostGIS and
        GDAL on your machine.
  \item Every installation error you met, copied in full, and what fixed
        it. If nothing went wrong, say so.
  \item The area table from Part 8, if you completed it.
  \item Your answers to the questions below.
\end{enumerate}

\subsection{Questions}

\begin{enumerate}
  \item In Labs 2 and 3, every visitor downloaded the whole GeoJSON file.
        What can a page backed by Django and PostGIS do instead? Give
        two examples from this lab.
  \item Why does \cd{settings.py} read the database password from an
        environment variable, instead of containing the password itself?
  \item In Part 4, the SQL ends with an index created
        \cd{USING GIST}. What does it do, and which query in this lab
        benefits from it?
  \item The loading script wraps some geometries in a
        \cd{MultiPolygon}. What would happen without that step, and why?
  \item Who should have an admin login for a district dataset like this,
        and who should not?
  \item Leaflet takes \cd{[lat, lng]} and \cd{Point} takes
        \cd{(lng, lat)}. Where in this lab is each order used, and how
        does the extension's URL avoid confusing them?
\end{enumerate}

\subsection{Marking}

\begin{tabularx}{\textwidth}{@{}X r@{}}
\toprule
\textbf{Item} & \textbf{Marks} \\
\midrule
Django project created and running                          & 10 \\
PostgreSQL and PostGIS installed, \cd{webgis} database ready & 15 \\
Django connected to PostGIS, migrations applied             & 15 \\
District model, loaded with 160 districts                   & 15 \\
Admin registered and explored, with the database edit       & 10 \\
Leaflet map drawn from the Django API                       & 15 \\
Report: versions, errors and answers to the six questions   & 20 \\
\midrule
Extension in Part 8                                         & +10 bonus \\
\bottomrule
\end{tabularx}

\vspace{1em}
Submit your repository link on LMS.


%==============================================================
\section*{Appendix A: Troubleshooting}
\addcontentsline{toc}{section}{Appendix A: Troubleshooting}

\begin{tabularx}{\textwidth}{@{}L{5cm} X@{}}
\toprule
\textbf{Message or symptom} & \textbf{Cause and fix} \\
\midrule
\cd{python} is not recognised &
  Python was installed without \textbf{Add python.exe to PATH}. Run the
  installer again, choose Modify, and tick it \\
Running scripts is disabled on this system &
  \cd{Set-ExecutionPolicy -Scope CurrentUser RemoteSigned}, then
  activate again \\
\cd{No module named django} &
  The virtual environment is not active. The prompt must start with
  \cd{(venv)} \\
Could not find the GDAL library &
  OSGeo4W is not in \cd{C:\textbackslash OSGeo4W}, or GDAL was not
  ticked in its installer. Check that
  \cd{C:\textbackslash OSGeo4W\textbackslash bin} contains a file named
  \cd{gdal} followed by digits and \cd{.dll} \\
Could not find the GEOS library &
  As above: \cd{geos\_c.dll} must be in the same \cd{bin} folder \\
\cd{fe\_sendauth: no password supplied} &
  \cd{WEBGIS\_DB\_PASSWORD} is not set in this window. Run \cd{setx}
  again if needed, then close and reopen VS Code \\
\cd{password authentication failed for user "postgres"} &
  The variable holds the wrong password. Run \cd{setx} with the right one
  and reopen VS Code \\
\cd{connection refused} on port 5432 &
  The PostgreSQL service is not running. Start it from Windows
  \textbf{Services} \\
\cd{type "geometry" does not exist} &
  \cd{CREATE EXTENSION postgis;} was not run in the \cd{webgis} database \\
\cd{database "webgis" does not exist} &
  Create it in pgAdmin, with exactly that name \\
\cd{TemplateDoesNotExist: districts/map.html} &
  The file must be at
  \cd{districts/templates/districts/map.html}, and \cd{districts} must be
  in \cd{INSTALLED\_APPS} \\
\cd{NoReverseMatch} for \cd{districts-geojson} &
  The \cd{name=} in \cd{urls.py} and the name in \cd{\{\% url \%\}} differ \\
Map page loads, no districts &
  Open \cd{/api/districts/} directly. An error there is a Python error;
  an empty collection means the loading script was not run \\
\bottomrule
\end{tabularx}

%==============================================================
\section*{Appendix B: Reference}
\addcontentsline{toc}{section}{Appendix B: Reference}

\subsection*{Commands}

\begin{code}
venv\Scripts\Activate.ps1          activate the environment
pip install -r requirements.txt    install what the project needs
python manage.py runserver         start the development server
python manage.py makemigrations    record a change to the models
python manage.py migrate           apply it to the database
python manage.py sqlmigrate APP N  show the SQL of a migration
python manage.py shell             a Python prompt with Django loaded
python manage.py createsuperuser   an admin account
\end{code}

\subsection*{The pieces and where they live}

\begin{tabularx}{\textwidth}{@{}L{5cm} X@{}}
\toprule
\textbf{File} & \textbf{Role} \\
\midrule
\cd{webgis/settings.py}  & Database, installed apps, OSGeo4W paths \\
\cd{webgis/urls.py}      & Sends \cd{admin/} to the admin, everything
                           else to the app \\
\cd{districts/models.py} & The table \\
\cd{districts/load.py}   & Fills it from GeoJSON \\
\cd{districts/admin.py}  & How the admin shows it \\
\cd{districts/views.py}  & What each URL returns \\
\cd{districts/urls.py}   & Which URL calls which view \\
\cd{districts/templates/} & The HTML Django sends to the browser \\
\bottomrule
\end{tabularx}

\subsection*{Optional, if you finish everything}

\begin{enumerate}
  \item Add a province drop-down to the map page that calls
        \cd{/api/districts/?province=...}, so the filtering happens in the
        database rather than in the browser.
  \item Add a \cd{Facility} model with a \cd{PointField}, register it in
        the admin, and add three facilities by clicking on the admin map.
  \item In the admin, create a second user who can change districts but
        not delete them, and log in as that user to check.
  \item Show the 25 districts without census data in grey, using the
        \cd{pop\_2023} property the API already sends.
\end{enumerate}

\end{document}
